\section{AMC 12 12B 2008}
        \begin{enumerate}
        	\item (\textit{AMC 12 12B 2008}) A basketball player made $5$ baskets during a game. Each basket was worth either $2$ or $3$ points. How many different numbers could represent the total points scored by the player?


 $\textbf{(A)}\ 2 \qquad \textbf{(B)}\ 3 \qquad \textbf{(C)}\ 4 \qquad \textbf{(D)}\ 5 \qquad \textbf{(E)}\ 6$

 \textbf{Answer: }E



\textbf{Solution 1}

If the basketball player makes $x$ three-point shots and $5-x$ two-point shots, he scores $3x+2(5-x)=10+x$ points. Clearly every value of $x$ yields a different number of total points. Since he can make any number of three-point shots between $0$ and $5$ inclusive, the number of different point totals is $6 \Rightarrow E$.

 



\textbf{Solution 2}

Stars and bars can also be utilized to solve this problem. Since we need to decide what number of 2's and 3's are scored, and there are a total of 5 shots. It can be written like such: \_ \_ \_ | \_ \_. Solving this, we get $6 \Rightarrow E$.

 


	\item (\textit{AMC 12 12B 2008}) A $4\times 4$ block of calendar dates is shown. The order of the numbers in the second row is to be reversed. Then the order of the numbers in the fourth row is to be reversed. Finally, the numbers on each diagonal are to be added. What will be the positive difference between the two diagonal sums?


 $\begin{tabular}[t]{|c|c|c|c|} \multicolumn{4}{c}{}\\\hline 1&2&3&4\\\hline 8&9&10&11\\\hline 15&16&17&18\\\hline 22&23&24&25\\\hline \end{tabular}$
$\textbf{(A)}\ 2 \qquad \textbf{(B)}\ 4 \qquad \textbf{(C)}\ 6 \qquad \textbf{(D)}\ 8 \qquad \textbf{(E)}\ 10$

 \textbf{Answer: }B



\textbf{Solution 1}

After reversing the numbers on the second and fourth rows, the block will look like this: 

 $\begin{tabular}[t]{|c|c|c|c|} \multicolumn{4}{c}{}\\\hline 1&2&3&4\\\hline 11&10&9&8\\\hline 15&16&17&18\\\hline 25&24&23&22\\\hline \end{tabular}$

 The positive difference between the two diagonal sums is then $(4+9+16+25)-(1+10+17+22)=3-1-1+3=\boxed{\textbf{(B) } 4}$.

 



\textbf{Solution 2}

Notice that at baseline the diagonals sum to the same number ($52$).  Therefore we need only compute the effect of the swap.  The positive difference between $9$ and $10$ is $1$ and the positive difference between $22$ and $25$ is $3$. Adding gives $1+3=\boxed{\textbf{(B) } 4}$



	\item (\textit{AMC 12 12B 2008}) A semipro baseball league has teams with $21$ players each. League rules state that a player must be paid at least $15,000$ dollars, and that the total of all players' salaries for each team cannot exceed $700,000$ dollars. What is the maximum possible salary, in dollars, for a single player?


 $\textbf{(A)}\ 270,000 \qquad \textbf{(B)}\ 385,000 \qquad \textbf{(C)}\ 400,000 \qquad \textbf{(D)}\ 430,000 \qquad \textbf{(E)}\ 700,000$

 \textbf{Answer: }C



\textbf{Solution}

We want to find the maximum any player could make, so assume that everyone else makes the minimum possible and that the combined salaries total the maximum of $700,000$
$700,000 = 20 * 15,000 + x$
$x = 400,000$

 The maximum any player could make is $400,000$ dollars $\Rightarrow C$



	\item (\textit{AMC 12 12B 2008}) On circle $O$, points $C$ and $D$ are on the same side of diameter $\overline{AB}$, $\angle AOC = 30^\circ$, and $\angle DOB = 45^\circ$. What is the ratio of the area of the smaller sector $COD$ to the area of the circle?


 
                    \begin{figure}[H]
                        \centering
                        \includegraphics[width=0.5\textwidth]{Problem/AMC_12_12B_2008_4_Problem_Figure_1.png}
                    \end{figure}
                    
$\textbf{(A)}\ \frac {2}{9} \qquad \textbf{(B)}\ \frac {1}{4} \qquad \textbf{(C)}\ \frac {5}{18} \qquad \textbf{(D)}\ \frac {7}{24} \qquad \textbf{(E)}\ \frac {3}{10}$

 \textbf{Answer: }D



\textbf{Solution}

$\angle COD = \angle AOB - \angle AOC - \angle BOD = 180^\circ - 30^\circ - 45^\circ = 105^\circ$. 

 Since a circle has $360^\circ$, the desired ratio is $\frac{105^\circ}{360^\circ}=\frac{7}{24} \Rightarrow D$. 

 


	\item (\textit{AMC 12 12B 2008}) A class collects $50$ dollars to buy flowers for a classmate who is in the hospital. Roses cost $3$ dollars each, and carnations cost $2$ dollars each. No other flowers are to be used. How many different bouquets could be purchased for exactly $50$ dollars?


 $\textbf{(A)}\ 1 \qquad \textbf{(B)}\ 7 \qquad \textbf{(C)}\ 9 \qquad \textbf{(D)}\ 16 \qquad \textbf{(E)}\ 17$

 \textbf{Answer: }C



\textbf{Solution}

The class could send $25$ carnations and no roses, $22$ carnations and $2$ roses, $19$ carnations and $4$ roses, and so on, down to $1$ carnation and $16$ roses. There are 9 total possibilities (from 0 to 16 roses, incrementing by 2 at each step), $\Rightarrow \boxed{C}$



	\item (\textit{AMC 12 12B 2008}) Postman Pete has a pedometer to count his steps. The pedometer records up to $99999$ steps, then flips over to $00000$ on the next step. Pete plans to determine his mileage for a year. On January $1$ Pete sets the pedometer to $00000$. During the year, the pedometer flips from $99999$ to $00000$ forty-four times. On December $31$ the pedometer reads $50000$. Pete takes $1800$ steps per mile. Which of the following is closest to the number of miles Pete walked during the year?


 $\textbf{(A)}\ 2500 \qquad \textbf{(B)}\ 3000 \qquad \textbf{(C)}\ 3500 \qquad \textbf{(D)}\ 4000 \qquad \textbf{(E)}\ 4500$

 \textbf{Answer: }A



\textbf{Solution}

Every time the pedometer flips, Pete has walked $100,000$ steps. Therefore, Pete has walked a total of $100,000 \cdot 44 + 50,000 = 4,450,000$ steps, which is $4,450,000/1,800 = 2472.2$ miles, which is the closest to the answer choice $\boxed{A}$.

 


	\item (\textit{AMC 12 12B 2008}) For real numbers $a$ and $b$, define $a$$$b = (a - b)\textasciicircum{}2$. What is $(x - y)\textasciicircum{}2$$$(y - x)^2$?


 $\textbf{(A)}\ 0 \qquad \textbf{(B)}\ x^2 + y^2 \qquad \textbf{(C)}\ 2x^2 \qquad \textbf{(D)}\ 2y^2 \qquad \textbf{(E)}\ 4xy$

 \textbf{Answer: }A



\textbf{Solution 1}

$\left[ (x-y)^2 - (y-x)^2 \right]^2$
$\left[ (x-y)^2 - (x-y)^2 \right]^2$
$[0]^2$
$0 \Rightarrow \textbf{(A)}$



	\item (\textit{AMC 12 12B 2008}) Points $B$ and $C$ lie on $\overline{AD}$. The length of $\overline{AB}$ is $4$ times the length of $\overline{BD}$, and the length of $\overline{AC}$ is $9$ times the length of $\overline{CD}$. The length of $\overline{BC}$ is what fraction of the length of $\overline{AD}$?


 $\textbf{(A)}\ \frac {1}{36} \qquad \textbf{(B)}\ \frac {1}{13} \qquad \textbf{(C)}\ \frac {1}{10} \qquad \textbf{(D)}\ \frac {5}{36} \qquad \textbf{(E)}\ \frac {1}{5}$

 \textbf{Answer: }C



\textbf{Solution}

Since $\overline{AB}=4\overline{BD}$ and $\overline{AB}+\overline{BD}=\overline{AD}$, $\overline{AB}=\frac{4}{5}\overline{AD}$.  

 Since $\overline{AC}=9\overline{CD}$ and $\overline{AC}+\overline{CD}=\overline{AD}$,  $\overline{AC}=\frac{9}{10}\overline{AD}$. 

 Thus, $\overline{BC}=\overline{AC}-\overline{AB}=\left(\frac{9}{10}-\frac{4}{5}\right)\overline{AD} = \frac {1}{10}\overline{AD} \Rightarrow C$. 

 


	\item (\textit{AMC 12 12B 2008}) Points $A$ and $B$ are on a circle of radius $5$ and $AB = 6$. Point $C$ is the midpoint of the minor arc $AB$. What is the length of the line segment $AC$?


 $\textbf{(A)}\ \sqrt {10} \qquad \textbf{(B)}\ \frac {7}{2} \qquad \textbf{(C)}\ \sqrt {14} \qquad \textbf{(D)}\ \sqrt {15} \qquad \textbf{(E)}\ 4$

 \textbf{Answer: }A



\textbf{Solution 1}

Let $\alpha$ be the angle that subtends the arc $AB$. By the law of cosines,$6^2=5^2+5^2-2\cdot 5\cdot 5\cos(\alpha)$ implies $\cos(\alpha) = 7/25$.

 The \href{https://artofproblemsolving.com/wiki/index.php?title=Trigonometric\_identities\#Half\_Angle\_Identities}{ half-angle formula} says that$\cos(\alpha/2) = \frac{\sqrt{1+\cos(\alpha)}}{2} = \sqrt{\frac{32/25}{2}} = \sqrt{\frac{16}{25}} = \frac{4}{5}$. The law of cosines tells us $AC = \sqrt{5^2+5^2-2\cdot 5\cdot 5\cdot \frac{4}{5}} = \sqrt{50-50\frac{4}{5}} = \sqrt{10}$, which is answer choice $\boxed{\text{A}}$.

 



\textbf{Solution 2}

Define $D$ as the midpoint of line segment $\overline{AB}$, and $O$ the center of the circle. Then $O$, $C$, and $D$ are collinear, and since $D$ is the midpoint of $AB$, $m\angle ODA=90\deg$ and so $OD=\sqrt{5^2-3^2}=4$. Since $OD=4$, $CD=5-4=1$, and so $AC=\sqrt{3^2+1^2}=\sqrt{10} \rightarrow \boxed{\text{A}}$.

 
                    \begin{figure}[H]
                        \centering
                        \includegraphics[width=0.5\textwidth]{Solution/AMC_12_12B_2008_9_Solution_Figure_1.png}
                    \end{figure}
                    



	\item (\textit{AMC 12 12B 2008}) Bricklayer Brenda would take $9$ hours to build a chimney alone, and bricklayer Brandon would take $10$ hours to build it alone. When they work together they talk a lot, and their combined output is decreased by $10$ bricks per hour. Working together, they build the chimney in $5$ hours. How many bricks are in the chimney?


 $\textbf{(A)}\ 500 \qquad \textbf{(B)}\ 900 \qquad \textbf{(C)}\ 950 \qquad \textbf{(D)}\ 1000 \qquad \textbf{(E)}\ 1900$

 \textbf{Answer: }B



\textbf{Solution}

Let $h$ be the number of bricks in the chimney. 

 Without talking, Brenda and Brandon lay $\frac{h}{9}$ and $\frac{h}{10}$ bricks per hour respectively, so together they lay $\frac{h}{9}+\frac{h}{10}-10$ per hour together. 

 Since they finish the chimney in $5$ hours, $h=5\left( \frac{h}{9}+\frac{h}{10}-10 \right)$. Thus, $h=900 \Rightarrow B$.

 


	\item (\textit{AMC 12 12B 2008}) A cone-shaped mountain has its base on the ocean floor and has a height of 8000 feet. The top $\frac{1}{8}$ of the volume of the mountain is above water. What is the depth of the ocean at the base of the mountain in feet?


 $\textbf{(A)}\ 4000 \qquad \textbf{(B)}\ 2000(4-\sqrt{2}) \qquad \textbf{(C)}\ 6000 \qquad \textbf{(D)}\ 6400 \qquad \textbf{(E)}\ 7000$

 \textbf{Answer: }A



\textbf{Solution}

In a cone, radius and height each vary inversely with increasing height (i.e. the radius of the cone formed by cutting off the mountain at $4,000$ feet is half that of the original mountain). Therefore, volume varies as the inverse cube of increasing height (expressed as a percentage of the total height of cone):$V_I\times \text{Height}^3 = V_N$

 Plugging in our given condition, $\frac{1}{8} = \text{Height}^3 \Rightarrow \text{Height} = \frac{1}{2}$.  

 $8000\cdot\frac{1}{2}=4000 \Rightarrow \boxed{\textbf{A}}$.

 \textbf{Faster Solution} The volume of the cone above water is $\frac 1 8$ that of the entire cone (mountain). These cones are obviously similar so the radius and height of the small cone must be $\sqrt[3]{\frac1 8}= \frac 1 2$ that of the large one. Because the height of the large cone is $8000$ the height of the small cone is $4000$. Thus the depth of the water is $8000-4000 = 4000 \Rightarrow\boxed{\text{A}}$ - AOPqghj

 


	\item (\textit{AMC 12 12B 2008}) For each positive integer $n$, the mean of the first $n$ terms of a sequence is $n$. What is the $2008$th term of the sequence?


 $\textbf{(A)}\ 2008 \qquad \textbf{(B)}\ 4015 \qquad \textbf{(C)}\ 4016 \qquad \textbf{(D)}\ 4030056 \qquad \textbf{(E)}\ 4032064$

 \textbf{Answer: }B



\textbf{Solution 1}

Letting $S_n$ be the nth partial sum of the sequence:

 $\frac{S_n}{n} = n$
$S_n = n^2$

 The only possible sequence with this result is the sequence of odd integers.

 $a_n = 2n - 1$
$a_{2008} = 2(2008) - 1 = 4015 \Rightarrow \textbf{(B)}$




\textbf{Solution 2}

Letting the sum of the sequence equal $a_1+a_2+\cdots+a_n$ yields the following two equations:

 $\frac{a_1+a_2+\cdots+a_{2008}}{2008}=2008$ and

 $\frac{a_1+a_2+\cdots+a_{2007}}{2007}=2007$.

 Therefore:

 $a_1+a_2+\cdots+a_{2008}=2008^2$ and $a_1+a_2+\cdots+a_{2007}=2007^2$

 Hence, by substitution, $a_{2008}=2008^2-2007^2=(2008+2007)(2008-2007)=4015(1)=4015\implies\boxed{\textbf{B}}$




\textbf{Solution 3}

Since the mean will be the sum of the first $n$ terms divided by $n$, and said mean also equals $n$, we know that the sum must be $n^2$. This means the sequence must compute squares, and this is done by the odd integer sequence $1+3+5+7+\ldots$. Therefore, we must find the 2008th odd number, which is found by $2n-1 \implies 2(2008)-1 = 4015$, so the answer is $\boxed{\textbf{B}}$~stress-couture

 


	\item (\textit{AMC 12 12B 2008}) Vertex $E$ of equilateral triangle $\triangle ABE$ is in the interior of unit square $ABCD$. Let $R$ be the region consisting of all points inside $ABCD$ and outside $\triangle ABE$ whose distance from $\overline{AD}$ is between $\frac{1}{3}$ and $\frac{2}{3}$. What is the area of $R$?


 $\textbf{(A)}\ \frac{12-5\sqrt{3}}{72} \qquad \textbf{(B)}\ \frac{12-5\sqrt{3}}{36} \qquad \textbf{(C)}\ \frac{\sqrt{3}}{18} \qquad \textbf{(D)}\ \frac{3-\sqrt{3}}{9} \qquad \textbf{(E)}\ \frac{\sqrt{3}}{12}$

 \textbf{Answer: }B



\textbf{Solution}

The region is the shaded area:

 
                    \begin{figure}[H]
                        \centering
                        \includegraphics[width=0.5\textwidth]{Solution/AMC_12_12B_2008_13_Solution_Figure_1.png}
                    \end{figure}
                    

 We can find the area of the shaded region by subtracting the pentagon from the middle third of the square.  The area of the middle third of the square is $\left(\frac13\right)(1)=\frac13$.  The pentagon can be split into a rectangle and an equilateral triangle.

 The base of the equilateral triangle is $\frac13$ and the height is $\left(\frac13\right)\left(\frac12\right)(\sqrt{3})=\frac{\sqrt{3}}{6}$.  Thus, the area is $\left(\frac{\sqrt3}{6}\right)\left(\frac13\right)\left(\frac12\right)=\frac{\sqrt3}{36}$.

 The base of the rectangle is $\frac13$ and the height is the height of the equilateral triangle minus the height of the smaller equilateral triangle.  This is:$\frac{\sqrt3}{2}-\frac{\sqrt3}{6}=\frac{\sqrt3}{3}$Therefore, the area of the shaded region is

 $\frac13-\frac{\sqrt3}{9}-\frac{\sqrt3}{36}=\boxed{\text{(B) }\frac{12-5\sqrt3}{36}}.$



	\item (\textit{AMC 12 12B 2008}) A circle has a radius of $\log_{10}{(a^2)}$ and a circumference of $\log_{10}{(b^4)}$. What is $\log_{a}{b}$?


 $\textbf{(A)}\ \frac{1}{4\pi} \qquad \textbf{(B)}\ \frac{1}{\pi} \qquad \textbf{(C)}\ \pi \qquad \textbf{(D)}\ 2\pi \qquad \textbf{(E)}\ 10^{2\pi}$

 \textbf{Answer: }C



\textbf{Solution}

Let $C$ be the circumference of the circle, and let $r$ be the radius of the circle. 

 Using log properties, $C=\log_{10}{(b^4)}=4\log_{10}{(b)}$ and $r=\log_{10}{(a^2)}=2\log_{10}{(a)}$. 

 Since $C=2\pi r$, $4\log_{10}{(b)}=2\pi\cdot2\log_{10}{(a)} \Rightarrow \log_{a}{b} = \frac{\log_{10}{(b)}}{\log_{10}{(a)}}=\pi \Rightarrow C$. 

 


	\item (\textit{AMC 12 12B 2008}) On each side of a unit square, an equilateral triangle of side length 1 is constructed. On each new side of each equilateral triangle, another equilateral triangle of side length 1 is constructed. The interiors of the square and the 12 triangles have no points in common. Let $R$ be the region formed by the union of the square and all the triangles, and $S$ be the smallest convex polygon that contains $R$. What is the area of the region that is inside $S$ but outside $R$?


 $\textbf{(A)} \; \frac {1}{4} \qquad \textbf{(B)} \; \frac {\sqrt {2}}{4} \qquad \textbf{(C)} \; 1 \qquad \textbf{(D)} \; \sqrt {3} \qquad \textbf{(E)} \; 2 \sqrt {3}$

 \textbf{Answer: }C



\textbf{Solution 1}


                    \begin{figure}[H]
                        \centering
                        \includegraphics[width=0.5\textwidth]{Solution/AMC_12_12B_2008_15_Solution_Figure_1.png}
                    \end{figure}
                    

 The equilateral triangles form trapezoids with side lengths $1, 1, 1, 2$ (half a unit hexagon) on each face of the unit square.  The four triangles "in between" these trapezoids are isosceles triangles with two side lengths $1$ and an angle of $30^{\circ}$ in between them, so the total area of these triangles (which is the area of $S - R$) is, $4 \left( \frac {1}{2} \sin 30^{\circ} \right) = 1$ which makes the answer $\boxed{C}$.

 



\textbf{Solution 2}


                    \begin{figure}[H]
                        \centering
                        \includegraphics[width=0.5\textwidth]{Solution/AMC_12_12B_2008_15_Solution_Figure_2.png}
                    \end{figure}
                    

 The area of the largest square is $(1+\sqrt{3})^2=4+2\sqrt{3}$. The area of region $R$ is $1+12\frac{1}{2}\frac{\sqrt{3}}{2}=1+3\sqrt{3}$. The total area of four small 45-45-90 triangles at corner is $4*\frac{1}{2}(\frac{\sqrt{3}+1-2}{2})^2=2-\sqrt{3}$. $S=4+2\sqrt{3}-(2-\sqrt{3})=2+3\sqrt{3}$, $S-R=1$.

 


	\item (\textit{AMC 12 12B 2008}) A rectangular floor measures $a$ by $b$ feet, where $a$ and $b$ are positive integers with $b > a$. An artist paints a rectangle on the floor with the sides of the rectangle parallel to the sides of the floor. The unpainted part of the floor forms a border of width $1$ foot around the painted rectangle and occupies half of the area of the entire floor. How many possibilities are there for the ordered pair $(a,b)$?


 $\textbf{(A)}\ 1\qquad\textbf{(B)}\ 2\qquad\textbf{(C)}\ 3\qquad\textbf{(D)}\ 4\qquad\textbf{(E)}\ 5$

 \textbf{Answer: }B



\textbf{Solution}

$A_{outer}=ab$
$A_{inner}=(a-2)(b-2)$
$A_{outer}=2A_{inner}$
$ab=2(a-2)(b-2)=2ab-4a-4b+8$
$0=ab-4a-4b+8$

 By Simon's Favorite Factoring Trick:

 $8=ab-4a-4b+16=(a-4)(b-4)$

 Since $8=1\times8$ and $8=2\times4$ are the only positive factorings of $8$.

 $(a,b)=(5,12)$ or $(a,b)=(6,8)$ yielding $\Rightarrow\textbf{(B)}$ $2$ solutions. Notice that because $b>a$, the reversed pairs are invalid.

 


	\item (\textit{AMC 12 12B 2008}) Let $A$, $B$ and $C$ be three distinct points on the graph of $y=x^2$ such that line $AB$ is parallel to the $x$-axis and $\triangle ABC$ is a right triangle with area $2008$. What is the sum of the digits of the $y$-coordinate of $C$?


 $\textbf{(A)}\ 16\qquad\textbf{(B)}\ 17\qquad\textbf{(C)}\ 18\qquad\textbf{(D)}\ 19\qquad\textbf{(E)}\ 20$

 \textbf{Answer: }C



\textbf{Solution}

Supposing $\angle A=90^\circ$, $AC$ is perpendicular to $AB$ and, it follows, to the $x$-axis, making $AC$  a segment of the line $x=m$. But that would mean that the coordinates of $C$ are $(m, m^2)$, contradicting the given that points $A$ and $C$ are distinct. So $\angle A$ is not $90^\circ$. By a similar logic, neither is $\angle B$. 

 This means that $\angle C=90^\circ$ and $AC$ is perpendicular to $BC$. Let C be the point $(n, n^2)$.  So the slope of $BC$ is the negative reciprocal of the slope of $AC$, yielding $m+n=\frac{1}{m-n}$ $\Rightarrow$ $m^2-n^2=1$. 

 Because $m^2-n^2$ is the length of the altitude of triangle $ABC$ from $AB$, and $2m$ is the length of $AB$, the area of $\triangle ABC=m(m^2-n^2)=2008$. Since $m^2-n^2=1$, $m=2008$. Substituting, $2008^2-n^2=1$ $\Rightarrow$ $n^2=2008^2-1=(2000+8)^2-1=4000000+32000+64-1=4032063$, whose digits sum to $18 \Rightarrow \textbf{(C)}$.

 


	\item (\textit{AMC 12 12B 2008}) A pyramid has a square base $ABCD$ and vertex $E$.  The area of square $ABCD$ is $196$, and the areas of $\triangle ABE$ and $\triangle CDE$ are $105$ and $91$, respectively.  What is the volume of the pyramid?


 $\textbf{(A)}\ 392 \qquad \textbf{(B)}\ 196\sqrt {6} \qquad \textbf{(C)}\ 392\sqrt {2} \qquad \textbf{(D)}\ 392\sqrt {3} \qquad \textbf{(E)}\ 784$

 \textbf{Answer: }E



\textbf{Solution}

Let $h$ be the height of the pyramid and $a$ be the distance from $h$ to $CD$. The side length of the base is $14$. The heights of $\triangle ABE$ and $\triangle CDE$ are $2\cdot105\div14=15$ and $2\cdot91\div14=13$, respectively. Consider a side view of the pyramid from $\triangle BCE$.  We have a systems of equations through the Pythagorean Theorem:

 $13^2-(14-a)^2=h^2 \\ 15^2-a^2=h^2$

 Setting them equal to each other and simplifying gives$-27+28a=225 \implies a=9$.

 Therefore, $h=\sqrt{15^2-9^2}=12$, and the volume of the pyramid is $\frac{bh}{3}=\frac{12\cdot 196}{3}=\boxed{784 \Rightarrow E}$.

 


	\item (\textit{AMC 12 12B 2008}) A function $f$ is defined by $f(z) = (4 + i) z^2 + \alpha z + \gamma$ for all complex numbers $z$, where $\alpha$ and $\gamma$ are complex numbers and $i^2 = - 1$. Suppose that $f(1)$ and $f(i)$ are both real. What is the smallest possible value of $| \alpha | + |\gamma |$?


 $\textbf{(A)} \; 1 \qquad \textbf{(B)} \; \sqrt {2} \qquad \textbf{(C)} \; 2 \qquad \textbf{(D)} \; 2 \sqrt {2} \qquad \textbf{(E)} \; 4 \qquad$

 \textbf{Answer: }B



\textbf{Solution 1:}

We need only concern ourselves with the imaginary portions of $f(1)$ and $f(i)$ (both of which must be 0). These are:

 
 \begin{align*} \text{Im}(f(1)) & = i+i\text{Im}(\alpha)+i\text{Im}(\gamma) \\ \text{Im}(f(i)) & = -i+i\text{Re}(\alpha)+i\text{Im}(\gamma) \end{align*}


 Let $p=\text{Im}(\gamma)$ and $q=\text{Re}{(\gamma)},$ then we know $\text{Im}(\alpha)=-p-1$ and $\text{Re}(\alpha)=1-p.$ Therefore 
 \[|\alpha|+|\gamma|=\sqrt{(1-p)^2+(-1-p)^2}+\sqrt{q^2+p^2}=\sqrt{2p^2+2}+\sqrt{p^2+q^2},\] which reaches its minimum $\sqrt 2$ when $p=q=0$ by the Trivial Inequality. Thus, the answer is $\boxed B.$




\textbf{Solution 2:}

$f(1)=4+i+\alpha+\gamma$
$f(i)=-4-i+\alpha \cdot i +\gamma$

 Since $f(1)$ and $f(i)$ are both real we get, 
 \[\alpha+\gamma=-i\]
 \[\alpha \cdot i+\gamma=i\]

 Solving, we get $\alpha=1-i$, $\gamma$ can be anything, to minimize the value we set $\gamma=0$, so then the answer is $\sqrt{1^2+1^2}=\sqrt{2}$. Thus, the answer is $\boxed{B}$

 By: Quaratinium

 


	\item (\textit{AMC 12 12B 2008}) Michael walks at the rate of $5$ feet per second on a long straight path. Trash pails are located every $200$ feet along the path. A garbage truck travels at $10$ feet per second in the same direction as Michael and stops for $30$ seconds at each pail. As Michael passes a pail, he notices the truck ahead of him just leaving the next pail. How many times will Michael and the truck meet?


 $\textbf{(A)}\ 4\qquad \textbf{(B)}\ 5\qquad \textbf{(C)}\ 6\qquad \textbf{(D)}\ 7\qquad \textbf{(E)}\ 8$

 \textbf{Answer: }B



\textbf{Solution 1}

Pick a coordinate system where Michael's starting pail is $0$ and the one where the truck starts is $200$.Let $M(t)$ and $T(t)$ be the coordinates of Michael and the truck after $t$ seconds.Let $D(t)=T(t)-M(t)$ be their (signed) distance after $t$ seconds.Meetings occur whenever $D(t)=0$.We have $D(0)=200$. 

 The truck always moves for $20$ seconds, then stands still for $30$. During the first $20$ seconds of the cycle the truck moves by $200$ feet and Michael by $100$, hence during the first $20$ seconds of the cycle $D(t)$ increases by $100$.During the remaining $30$ seconds $D(t)$ decreases by $150$.

 From this observation it is obvious that after four full cycles, i.e. at $t=200$, we will have $D(t)=0$ for the first time.

 During the fifth cycle, $D(t)$ will first grow from $0$ to $100$, then fall from $100$ to $-50$. Hence Michael overtakes the truck while it is standing at the pail. 

 During the sixth cycle, $D(t)$ will first grow from $-50$ to $50$, then fall from $50$ to $-100$. Hence the truck starts moving, overtakes Michael on their way to the next pail, and then Michael overtakes the truck while it is standing at the pail. 

 During the seventh cycle, $D(t)$ will first grow from $-100$ to $0$, then fall from $0$ to $-150$. Hence the truck meets Michael at the moment when it arrives to the next pail.

 Obviously, from this point on $D(t)$ will always be negative, meaning that Michael is already too far ahead. Hence we found all $\boxed{5 \Longrightarrow B}$ meetings.

 The movement of Michael and the truck is plotted below: Michael in blue, the truck in red.

 
                    \begin{figure}[H]
                        \centering
                        \includegraphics[width=0.5\textwidth]{Solution/AMC_12_12B_2008_20_Solution_Figure_1.png}
                    \end{figure}
                    




\textbf{Solution 2}

The truck takes $20$ seconds to go from one pail to the next and then stops for $30$ seconds at the new pail. Thus it sets off from a pail every 50 sec. Let $t$ denote the time elapsed and write $t=50k + \Delta$, where $\Delta \in [0,50)$. In  this time Michael has traveled $5t = 250k+5\Delta$ feet. What about the truck? In the first $50k$ seconds the truck covers $k$ pails, i.e. $200k$ feet so it moves $200+200k$ feet from Michael's starting point. Then we have two cases: 

 \textbf{(a)} if $\Delta < 20$, then the truck travels an additional $10\Delta$ feet. For them to intersect we must have $200+200k + 10\Delta = 250k + 5\Delta$. Solving, we get $\Delta = 10k - 40$. Since $\Delta$ must lie in the interval $[0,20)$ we get $k \in \{4, 5\}$.

 \textbf{(b)} if $\Delta \in [20, 50)$, then the truck travels an additional $200$ feet. For them to intersect we must have $200+200k + 200 = 250k + 5\Delta$. Solving, we get $\Delta = 10(8-k)$. Since $\Delta$ must lie in the interval $[20,50)$ we get $k \in \{4, 5, 6\}$. 

 Thus Michael intersects with the truck $5$ times, which is option $\boxed{\textbf{(B)}}$.

 



\textbf{Solution 3}

We make a chart by seconds in increments of ten. 

 $\begin{tabular}{c|c|c}Time (s) &Distance of Michael&Distance of Garbage\\ \hline 0&0&200\\ 10&50&300\\ 20&100&400\\ 30&150&400\\ 40&200&400\\ 50&250&400\\ 60&300&500\\ 70&350&600\\ 80&400&600\\ 90&450&600\\ 100&500&600\\ 110&550&700\\ 120&600&800\\ 130&650&800\\ 140&700&800\\ 150&750&800\\ 160&800&900\\ 170&850&1000\\ 180&900&1000\\ 190&950&1000\\ 200&1000&1000\\ 210&1050&1100\\ 220&1100&1200\\ 230&1150&1200\\ 240&1200&1200\\ 250&1250&1200\\ 260&1300&1300\\ 270&1350&1400\\ 280&1400&1400\\ 290&1450&1400\\ 300&1500&1400\\ 310&1550&1500\\ 320&1600&1600\\ 330&1650&1600\\ 340&1700&1600\\ 350&1750&1600\\ 360&1800&1700\\ 370&1850&1800\\ 380&1900&1800\\ 390&1950&1800\\ 400&2000&1800\\ \end{tabular}$

 
Notice at 200, 240, 260, 280, and 320 seconds, Michael and the garbage truck meet. It is clear that they met at these times, and will meet no more. Thus the answer is $\boxed{\textbf{(B)}}$.

 


 



\textbf{Solution 4 (Educated Guess)}

We can record the time of the truck and the time of Michael at each trash pail. 

 $\begin{tabular}{c|c|c}Distance (s) &Truck Time&Michael Time\\ \hline 400&20, 50&40\\ 600&70, 100&80\\ 800&120, 150&120\\ 1000&170, 200&160\\ 1200&220, 250&200\\ 1400&270, 300&240\\ 1600&320, 350&320\\ \end{tabular}$

 Since Micheal moves faster than the truck, the truck will not catch up to Michael afterwards.

 We see that there are 4 correct times and we add 1 because they also pass between pails once which is option $\boxed{\textbf{(B)}}$



	\item (\textit{AMC 12 12B 2008}) Two circles of radius 1 are to be constructed as follows. The center of circle $A$ is chosen uniformly and at random from the line segment joining $(0,0)$ and $(2,0)$. The center of circle $B$ is chosen uniformly and at random, and independently of the first choice, from the line segment joining $(0,1)$ to $(2,1)$. What is the probability that circles $A$ and $B$ intersect?


 $\textbf{(A)} \; \frac {2 + \sqrt {2}}{4} \qquad \textbf{(B)} \; \frac {3\sqrt {3} + 2}{8} \qquad \textbf{(C)} \; \frac {2 \sqrt {2} - 1}{2} \qquad \textbf{(D)} \; \frac {2 + \sqrt {3}}{4} \qquad \textbf{(E)} \; \frac {4 \sqrt {3} - 3}{4}$

 \textbf{Answer: }E



\textbf{Solution 1}

Circles centered at $A$ and $B$ will overlap if $A$ and $B$ are closer to each other than if the circles were tangent. The circles are tangent when the distance between their centers is equal to the sum of their radii. Thus, the distance from $A$ to $B$ will be $2$. Since $A$ and $B$ are separated by $1$ vertically, they must be separated by $\sqrt{3}$ horizontally. Thus, if $|A_x-B_x|<\sqrt{3}$, the circles intersect.

 Now, plot the two random variables $A_x$ and $B_x$ on the coordinate plane. Each variable ranges from $0$ to $2$. The circles intersect if the variables are within $\sqrt{3}$ of each other. Thus, the area in which the circles don't intersect is equal to the total area of two small triangles on opposite corners, each of area $\frac{(2-\sqrt{3})^2}{2}$. So, the total area of the 2 triangles sums to $(2-\sqrt{3})^2$. Since the total $2 \times 2$ square has an area of $4$, the probability of the circles not intersecting is $\frac{(2-\sqrt{3})^2}{4}$. But remember, we want the probability that they do intersect. We conclude the probability the circles intersect is:
 \[1-\frac{(2-\sqrt{3})^2}{4}=\boxed{\textbf{(E)}\frac{4\sqrt{3}-3}{4}}.\]




\textbf{Solution 2}

Two circles intersect if the distance between their centers is less than the sum of their radii. In this problem, $A$ and $B$ intersect if $\sqrt{(A_X-B_X)^2+(A_Y-B_Y)^2}\leq2 \Rightarrow (A_X-B_X)^2+1\leq4 \Rightarrow (A_X-B_X)\leq \sqrt{3}.$

 In other words, the two chosen $x$-coordinates must differ by no more than $\sqrt{3}$. To find this probability, we divide the problem into cases:

 1) $A_X$ is on the interval $(0,2-\sqrt{3})$. The probability that $B_X$ falls within the desired range for a given $A_X$ is $A_X$ (on the left) $+\sqrt{3}$ (on the right) all over $2$ (the range of possible values). The total probability for this range is the sum of all these probabilities of $B_X$ (over the range of $A_X$) divided by the total range of $A_X$ (which is $2$). Thus, the total probability for this interval is 
 \[\frac{1}{2}\left(\int_{0}^{2-\sqrt{3}} \frac{x+\sqrt{3}}{2}\,dx\right) =\frac{1}{2}\left(x^2/4+\frac{x\sqrt{3}}{2} \Big |_{0}^{2-\sqrt{3}}\right)\] 
 \[= \frac{1}{2}(\frac{4-4\sqrt{3}+3}{4} + \sqrt{3}-\frac{3}{2}) =\frac{1}{8}.\]2) $A_X$ is on the interval $(2-\sqrt{3},\sqrt{3})$. In this case, any value of $B_X$ will do, so the probability for the interval is simply $\frac{\sqrt{3}-(2-\sqrt{3})}{2} = \sqrt{3}-1$.

 3) $A_X$ is on the interval $(\sqrt{3},2)$. This is identical, by symmetry, to case 1.

 The total probability is therefore $\sqrt{3}-1+\frac{1}{4} = \frac{4\sqrt{3}-3}{4} \Rightarrow \boxed{E}$




\textbf{Solution 3}

We first calculate the probability that the circles at are their maximum possible distance while still intersecting. Since the difference in heights is 1, and the two radii add up to 2 (at their point of tangency), we can see that the maximum possible distance for the two centers is $\sqrt3$. If we look at possible placements for the lower point, we see that at placement $(0,0)$ there is exactly $\sqrt3$ of space wherein we can put the top point. Similarly, at point $(1,0)$ we can put the top point anywhere, resulting in $2$ of space. If one looks at a diagram (which I can't put here), we can see that the space available for the top point increases linearly from $(0,0)$ to $(2-\sqrt3,0)$. From then on to $(1,0)$ it stops at $2$ (this is symmetric for the other half of placements). If one draws a plot of the space available (which I again can't show here) with x determined by the x-value of the bottom point and y determined by space available for the top point, we notice two things: first, since for each (theoretically infinitesimally) small x-slice there is a line segment determining how much space there is available, if we sum up all of the line segments and divide by the total space (technically speaking by enumerating all the valid combinations and dividing by all of the general combinations), we will end up with a plot that looks like a rectangle with two small triangles cut out of it. Two, that since this is symmetric, we can just focus on the first half as the ratio of valid to total is the same. We see that the rectangle has a height of $2$ and a width of $1$, thus its area is 2. The triangle's height is $2-\sqrt3$, and its width is the same thus its area is $(2-\sqrt3)^2/2$ => $(7-4\sqrt3)/2$. Now we simply subtract to obtain $(4\sqrt3-3)/2$ is the valid area, thus $(4\sqrt3-3)/4 \Rightarrow \boxed{E}$ is our solution.

 


	\item (\textit{AMC 12 12B 2008}) A parking lot has 16 spaces in a row. Twelve cars arrive, each of which requires one parking space, and their drivers chose spaces at random from among the available spaces. Auntie Em then arrives in her SUV, which requires 2 adjacent spaces. What is the probability that she is able to park?


 $\textbf{(A)} \; \frac {11}{20} \qquad \textbf{(B)} \; \frac {4}{7} \qquad \textbf{(C)} \; \frac {81}{140} \qquad \textbf{(D)} \; \frac {3}{5} \qquad \textbf{(E)} \; \frac {17}{28}$

 \textbf{Answer: }E



\textbf{Solution}

Auntie Em won't be able to park only when none of the four available spots touch. We can form a \href{https://artofproblemsolving.com/wiki/index.php?title=Bijection}{bijection} between all such cases and the number of ways to pick four spots out of 13: since none of the spots touch, remove a spot from between each of the cars. From the other direction, given four spots out of 13, simply add a spot between each.So the probability she can park is 
 \[1-\frac{{13 \choose 4}}{{16 \choose 4}}=1-\frac{13\cdot12\cdot11\cdot10}{16\cdot15\cdot14\cdot13}=1-\frac{11}{28}={\textbf{(E)}\frac{17}{28}}.\]

 (Bijection: When elements of two sets are perfectly paired with each other, i.e. each and every element from both sets has exactly one match in the other, and no elements are left out. In the context of this problem, this means the number of distinct ways to order the cars such that no two spaces are adjacent is exactly the number of ways to pick 4 spots out of 13.)

 \textbf{Further Explanation} Instead of trying to park the cars one at a time(considering the number of choices for the first car, then the second, and so on), imagine that all of the cars are already parked, and ignore the empty parking spaces for now. This will take up $12$ parking spaces, illustrated as asterisks:

 
 \[************\]

 We know that 4 spaces must be left empty, and we care about the arrangements where no empty spaces are next to one another. Notice that this is simply the number of ways to put $4$|s into $13$ slots (including the "slot" in front of the first car as well):

 
 \[|*****|*|*|***** \text{(this works)}\]

 \[|******||*|***** \text{(this doesn't)}\]

 In other words, after a space between two cars has been selected to be empty, that space is no longer an option for the next selection. The number of available spaces will decrease by one every time, so we have $13*12*11*10$ ways to arrange the empty spaces. But the empty spaces are identical, so we must divide by the number of orders in which we can place the empty spaces, or $4!$. This is equivalent to $\binom{13}{4}$ total ways to arrange $4$ empty spaces into $13$ "slots". Then we continue similarly to the above.

 



\textbf{Solution 2 (distributions)}

This solution is pretty similar to the above. So there are 4 spots remaining, and it would be hard to count all the combinations where Auntie Em could park, but you can count all the combinations where Auntie Em can't park. Since the 12 cars are indistinguishable, we can use distributions here.

 There must be $a_1 \geq 0$ cars, then an empty spot, then $d_2 \geq 1$ cars, then an empty spot, then $d_3 \geq 1$ cars, then an empty spot, then $d_4 \geq 1$ cars, then an empty spot, then $a_5 \geq 0$ cars.

 
 \[a_1 + d_2 + d_3 + d_4 + a_5 = 12\]

 To remove all the restrictions, let $a_n = d_n - 1$.

 
 \[a_1 + a_2 + a_3 + a_4 + a_5 = 9\]

 We can now use stars and bars on these values (9 "stars", and 5 - 1 = 4 dividers), to get $\binom{13}{4}$ possibilities where she can't park. There are $\binom{16}{4}$ possibilities in total.

 
 \[\frac{\binom{13}{4}}{\binom{16}{4}} = \frac{11}{28}\]

 Subtracting that from 1 to get the probability she can park, the correct answer is $\boxed{E}$.

 


	\item (\textit{AMC 12 12B 2008}) The sum of the base-$10$ logarithms of the divisors of $10^n$ is $792$. What is $n$?


 $\textbf{(A)}\ 11\qquad \textbf{(B)}\ 12\qquad \textbf{(C)}\ 13\qquad \textbf{(D)}\ 14\qquad \textbf{(E)}\ 15$

 \textbf{Answer: }A



\textbf{Solution 1}

Every factor of $10^n$ will be of the form $2^a \times 5^b , a\leq n , b\leq n$. Not all of these base ten logarithms will be rational, but we can add them together in a certain way to make it rational. Recall the logarithmic property $\log(a \times b) = \log(a)+\log(b)$. For any factor $2^a \times 5^b$, there will be another factor $2^{n-a} \times 5^{n-b}$. Note this is not true if $10^n$ is a perfect square. When these are multiplied, they equal $2^{a+n-a} \times 5^{b+n-b} = 10^n$. $\log 10^n=n$ so the number of factors divided by 2 times n equals the sum of all the factors, 792. 

 There are $n+1$ choices for the exponent of 5 in each factor, and for each of those choices, there are $n+1$ factors (each corresponding to a different exponent of 2), yielding $(n+1)^2$ total factors. $\frac{(n+1)^2}{2}*n = 792$. We then plug in answer choices and arrive at the answer $\boxed {11}$




\textbf{Solution 2}

We are given 
 \[\log_{10}d_1 + \log_{10}d_2 + \ldots + \log_{10}d_k = 792\] The property $\log(ab) = \log(a)+\log(b)$ now gives 
 \[\log_{10}(d_1 d_2\cdot\ldots d_k) = 792\] The product of the divisors is (from elementary number theory) $a^{d(n)/2}$ where $d(n)$ is the number of divisors. Note that $10^n = 2^n\cdot 5^n$, so $d(n) = (n + 1)^2$. Substituting these values with $a = 10^n$ in our equation above, we get $n(n + 1)^2 = 1584$, from whence we immediately obtain $\framebox[1.2\width]{(A)}$ as the correct answer.

 



\textbf{Solution 3}

For every divisor $d$ of $10^n$, $d \le \sqrt{10^n}$, we have $\log d + \log \frac{10^n}{d} = \log 10^n = n$. There are $\left \lfloor \frac{(n+1)^2}{2} \right \rfloor$ divisors of $10^n = 2^n \times 5^n$ that are $\le \sqrt{10^n}$. After casework on the parity of $n$, we find that the answer is given by $n \times \frac{(n+1)^2}{2} = 792 \Longrightarrow n = 11\ \mathrm{(A)}$.

 



\textbf{Solution 4}

The sum is
 \[\sum_{p=0}^{n}\sum_{q=0}^{n} \log(2^p5^q) = \sum_{p=0}^{n}\sum_{q=0}^{n}(p\log(2)+q\log(5))\]
 \[= \sum_{p=0}^{n} ((n+1)p\log(2) + \frac{n(n+1)}{2}\log(5))\]
 \[= \frac{n(n+1)^2}{2} \log(2) + \frac{n(n+1)^2}{2} \log(5)\]
 \[= \frac{n(n+1)^2}{2}\]Trying for answer choices we get $n=11$




\textbf{Solution 5}

Let integer $d$ be the number of divisors $10^n$ has. Then, we set up $\frac{d}{2}$ pairs of divisors such that each pair $(a,b)$ satisfies $ab = 10^n$ -- ex. $(1, 10^n), (2, 5*10^{n-1})$, etc. Then the sum of the base-$10$ logarithms is $\sum_{j=1}^{\frac{d}{2}} \log_{10} a_j + \log_{10} b_j = \sum_{j=1}^{\frac{d}{2}} \log_{10} a_j b_j = \sum_{j=1}^{\frac{d}{2}} \log_{10} 10^n = \sum_{j=1}^{\frac{d}{2}} n = \frac{nd}{2}$We can use the property that only perfect squares have an odd number of factors, as for perfect square $p^2$, we have ordered pair $(p,p)$ that works. For even $n$, then, $10^{\frac{n}{2}}$ can be multiplied by itself to get $10^n$, so $d$ is odd. But, in our summation, $\frac{d}{2}$ does not exist for even $n$ as $d$ is then odd, so $d$ must be even. And, since $\frac{nd}{2}$ = $792$, We want to find a $d$ for our odd options $n$ such that $nd$ = $1584$. For $n=11$, $d=144$ works, and integer $d$ can not be found for other odd $n$. So, we get $\framebox[1.2\width]{(A) 11}$

 \textbf{Alternative thinking} After arriving at the equation $n(n+1)^2 = 1584$, notice that all of the answer choices are in the form $10+k$, where $k$ is $1-5$. We notice that the ones digit of $n(n+1)^2$ is $4$, and it is dependent on the ones digit of the answer choices. Trying $1-5$ for $n$, we see that only $1$ yields a ones digit of $4$, so our answer is $11$.

 


	\item (\textit{AMC 12 12B 2008}) Let $A_0 = (0,0)$. Distinct points $A_1,A_2,\ldots$ lie on the $x$-axis, and distinct points $B_1,B_2,\ldots$ lie on the graph of $y = \sqrt {x}$. For every positive integer $n$, $A_{n - 1}B_nA_n$ is an equilateral triangle. What is the least $n$ for which the length $A_0A_n\ge100$?


 $\textbf{(A)}\ 13\qquad \textbf{(B)}\ 15\qquad \textbf{(C)}\ 17\qquad \textbf{(D)}\ 19\qquad \textbf{(E)}\ 21$

 \textbf{Answer: }C



\textbf{Solution 1}

Let $a_n=|A_{n-1}A_n|$. We need to rewrite the recursion into something manageable. The two strange conditions, $B$'s lie on the graph of $y=\sqrt{x}$ and $A_{n-1}B_nA_n$ is an equilateral triangle, can be compacted as follows: 
 \[\left(a_n\frac{\sqrt{3}}{2}\right)^2=\frac{a_n}{2}+a_{n-1}+a_{n-2}+\cdots+a_1\]which uses $y^2=x$, where $x$ is the height of the equilateral triangle and therefore $\frac{\sqrt{3}}{2}$ times its base.

 The relation above holds for $n=k$ and for $n=k-1$ $(k>1)$, so 
 \[\left(a_k\frac{\sqrt{3}}{2}\right)^2-\left(a_{k-1}\frac{\sqrt{3}}{2}\right)^2=\]
 \[=\left(\frac{a_k}{2}+a_{k-1}+a_{k-2}+\cdots+a_1\right)-\left(\frac{a_{k-1}}{2}+a_{k-2}+a_{k-3}+\cdots+a_1\right)\]Or, 
 \[a_k-a_{k-1}=\frac23\] This implies that each segment of a successive triangle is $\frac23$ more than the last triangle. To find $a_{1}$, we merely have to plug in $k=1$ into the aforementioned recursion and we have $a_{1} - a_{0} = \frac23$. Knowing that $a_{0}$ is $0$, we can deduce that $a_{1} = 2/3$.Thus, $a_n=\frac{2n}{3}$, so $A_0A_n=a_n+a_{n-1}+\cdots+a_1=\frac{2}{3} \cdot \frac{n(n+1)}{2} = \frac{n(n+1)}{3}$. We want to find $n$ so that $n^2<300<(n+1)^2$. $n=\boxed{17}$ is our answer.

 



\textbf{Solution 2}

Consider two adjacent equilateral triangles obeying the problem statement. For each, drop an altitude to the $x$ axis and denote the resulting heights $h_n$ and $h_{n+1}$. From 30-60-90 rules, the distance between the points where these altitudes meet the x-axis is 
 \[\frac{h_{n+1}}{\sqrt{3}}+\frac{h_n}{\sqrt{3}} = \frac{h_{n+1}+h_n}{\sqrt{3}}\]

 But the square root curve means that this distance is also expressible as $h_{n+1}^2-h_n^2$ (the $x$ coordinates are the squares of the heights). Setting these expressions equal and dividing throughout by $h_{n+1}+h_n$ leaves $h_{n+1}-h_n=\frac{1}{\sqrt{3}}$. So the difference in height of successive triangles is $\frac{1}{\sqrt{3}}$, meaning their bases are wider by $2/3$ units each time. From here, one can proceed as in Solution 1 to arrive at $n=\boxed{17}$.

 



\textbf{Solution 3}

Note that $A_{1}$ is of the form $(2x,0)$ for some $x$, and thus $B_{1}$ is of the form $(x, x \sqrt{3}).$ Then, we are told that $B_{1}$ lies on the graph of $y = \sqrt{x}$, so 
 \[(x \sqrt{3})^{2} = x.\] Solving for x, we get that $x = \frac{1}{3},$ and so $A_{1} = (2/3,0)$. 

 Now, similarly to before, let $|A_{1}A_{2}|=2y.$ Then, $B_{2} = (y+\frac{2}{3},y \sqrt{3})$, and so 
 \[(y \sqrt{3})^{2} = (y+\frac{2}{3}).\]Solving using the quadratic formula gives 
 \[y = \frac{-(-1) \pm \sqrt{1^{2}-4(3)(-\frac{2}{3})}}{6} = \frac{1 \pm \sqrt{1^{2}+8}}{6} = \frac{2}{3}.\]Then, $2y = \frac{4}{3},$ so $A_{2} = (2,0).$

 In general, if $A_{n} = (a_{n},0)$ for all integers $n$, and  $|A_{n}A_{n+1}| = 2x$ for some real number $x$, we have the following equation for $x$:
 \[(x \sqrt{3})^{2} = x+a_{n},\]which give us when plugged into the quadratic formula gives
 \[x = \frac{1 \pm \sqrt{1+12a_{n}}}{6},\]but since $x$ must be positive, we have that
 \[x = \frac{1 + \sqrt{1+12a_{n}}}{6},\]and so 
 \[a_{n+1} = a_{n}+2x = a_{n}+\frac{1 + \sqrt{1+12a_{n}}}{3}.\]Computing a few terms of $a_{n}$ using this method gives $a_{3} = 4,$ $a_{4} = \frac{20}{3},$ and $a_{5} = 10$. 

 Notice how all our $a_{n}$ terms so far are rational, even though there is an abundance of radicals in the recurrence. This motivates us to look at our discriminants in the quadratic formula that is solved for $x$. 

 The discriminant of 
 \[(x \sqrt{3})^{2} = x+a_{2}\] is 
 \[1+12 \cdot \frac{2}{3} = 9.\] Similarly, the discriminant of 
 \[(x \sqrt{3})^{2} = x+a_{3}\] is $1+12 \cdot a_{3} = 25,$ and $1+12 \cdot a_{4} = 49$. 

 Note how our results keep coming out as the squares of the odd integers. Moreover, it seems that 
 \[1+12a_{n} = (2n+1)^{2}.\]We will prove this with induction. The base case, $n=2$, we have already verified.

 Now, for the Inductive step, assume that 
 \[1+12a_{n-1} = (2n-1)^{2}.\] for some integer $n-1$. We will prove that this is true for $n$ as well. 

 Plugging this into our recurrence formula gives us
 \[a_{n} = a_{n}+2x = a_{n-1}+\frac{1 + \sqrt{1+12a_{n-1}}}{3}\]
 \[a_{n} = \frac{(2n-1)^{2}-1}{12} + \frac{1 + (2n-1)}{3}\]
 \[a_{n} = \frac{4n^{2}-4n+1-1}{12} + \frac{8n}{12}\]
 \[a_{n} = \frac{4n^{2}+4n+1-1}{12}\]
 \[a_{n} = \frac{(2n+1)^{2}-1}{12}.\]Therefore, we have proved our claim. Now, we have that $|A_{0}A_{n}|=a_{n},$ so we just need the least integer $n$ so that 
 \[a_{n} > 100,\]or
 \[(2n+1)^{2} > 1201.\]Then, we see that $35^{2} = 1225$ is the smallest odd square larger than $1201.$ Therefore, we have $2n+1=35$, so $n = \boxed{17}.$

 -Mr.Sharkman

 



\textbf{Solution 4 (Pattern Observation)}

We can iteratively calculate out the first few $A_i$ and $B_i$. We know that $A_0 = (0,0)$ and the line through $A_0$ and $B_1$ needs to make a $60^{\circ}$ angle with the x-axis (because the triangle is equilateral). The equation of a line that makes an angle $\theta$ with the x-axis and passes through the origin has equation $\tan(\theta)x$, so the line through $A_0$ and $B_1$ has equation $\sqrt{3}x$. We set this equal to $\sqrt{x}$ (to find where the two curves intersect) and, when solving, find that $x = \frac{1}{3}$. Therefore, $B_1 = (\frac{1}{3}, \frac{\sqrt{3}}{3})$. By equilateral triangle properties, we must then have that $A_1 = (\frac{2}{3}, 0)$. To find $B_2$, we find the equation of the line through $A_1$ that makes a $60^{\circ}$ angle with the $x$-axis. This is the same line as the one through the origin (which we already found) shifted $\frac{2}{3}$ to the right, so it has equation $y = \sqrt{3}(x - \frac{2}{3})$. Setting this equal to $\sqrt{x}$ and solving for $x$, we get $x = \frac{4}{3}$, so $B_2 = (\frac{4}{3}, \frac{2\sqrt{3}}{3})$. By equilateral triangle properties, we have that $A_2 = (2, 0)$. Repeating this process, we find that $B_3 = (3, \sqrt{3})$ and $A_3 = (4,0)$. At this point, we notice that the $y$-coordinate of each $B_i$ is $\frac{\sqrt{3}}{3}$ more than that of $B_{i-1}$, so $B_4 = (\frac{16}{3}, \frac{4\sqrt{3}}{3})$ and $A_4 = (\frac{20}{3}, 0)$. It may be helpful to continue calculating out the $A_i$ and $B_i$ if the pattern in the $x$-coordinates of the $A_i$ isn't visible yet. At this point, it can be noticed that the differences between the x coordinates of consecutive $A_i$ form an arithmetic sequence ($x_{A_1} - x_{A_0} = \frac{2}{3}$, $x_{A_2} - x_{A_1} = \frac{4}{3}$, etc.). Additionally, the numerator of each of the fractions is a consecutive even number, and the sum of the first $n$ even numbers is $n(n+1)$. Thus, we must find the least integral solution to $\frac{n(n+1)}{3} \geq 100$. This least solution is $\boxed{\textbf{(C)}~17}$.

 


	\item (\textit{AMC 12 12B 2008}) Let $ABCD$ be a trapezoid with $AB||CD$, $AB = 11$, $BC = 5$, $CD = 19$, and $DA = 7$. Bisectors of $\angle A$ and $\angle D$ meet at $P$, and bisectors of $\angle B$ and $\angle C$ meet at $Q$. What is the area of hexagon $ABQCDP$?


 $\textbf{(A)}\ 28\sqrt {3}\qquad \textbf{(B)}\ 30\sqrt {3}\qquad \textbf{(C)}\ 32\sqrt {3}\qquad \textbf{(D)}\ 35\sqrt {3}\qquad \textbf{(E)}\ 36\sqrt {3}$

 \textbf{Answer: }B



\textbf{Solution}


                    \begin{figure}[H]
                        \centering
                        \includegraphics[width=0.5\textwidth]{Solution/AMC_12_12B_2008_25_Solution_Figure_1.jpg}
                    \end{figure}
                    

 Note: In the image AB and CD have been swapped from their given lengths in the problem. However, this doesn't affect any of the solving.

 Drop perpendiculars to $CD$ from $A$ and $B$, and call the intersections $X,Y$ respectively. Now, $DA^2-BC^2=(7-5)(7+5)=DX^2-CY^2$ and $DX+CY=19-11=8$. Thus, $DX-CY=3$.We conclude $DX=\frac{11}{2}$ and $CY=\frac{5}{2}$.To simplify things even more, notice that $90^{\circ}=\frac{\angle D+\angle A}{2}=180^{\circ}-\angle APD$, so $\angle P=\angle Q=90^{\circ}$.

 Also, 
 \[\sin(\angle PAD)=\sin(\frac12\angle XDA)=\sqrt{\frac{1-\cos(\angle XDA)}{2}}=\sqrt{\frac{3}{28}}\]So the area of $\triangle APD$ is: 
 \[R\cdot c\sin a\sin b =\frac{7\cdot7}{2}\sqrt{\frac{3}{28}}\sqrt{1-\frac{3}{28}}=\frac{35}{8}\sqrt{3}\]

 Over to the other side: $\triangle BCY$ is $30-60-90$, and is therefore congruent to $\triangle BCQ$. So $[BCQ]=\frac{5\cdot5\sqrt{3}}{8}$.

 The area of the hexagon is clearly 
 \begin{align*} [ABCD]-([BCQ]+[APD]) &=\frac{15\cdot5\sqrt{3}}{2}-\frac{60\sqrt{3}}{8}\\ &=30\sqrt{3}\implies\boxed{\mathrm{B}} \end{align*}


 Note: Once $DY$ is found, there is no need to do the trig. Notice that the hexagon consists of two trapezoids, $ABPQ$ and $CDPQ$. $PQ = \frac{19-7-5 +11}{2} = 9$. The height is one half of $BY$ which is $\frac{5\sqrt{3}}{4}$. So the area is 
 \[\frac{1}{2} \cdot \frac{5\sqrt{3}}{4}(19+9+11+9)=30\sqrt{3}\]




\textbf{Solution 2}


                    \begin{figure}[H]
                        \centering
                        \includegraphics[width=0.5\textwidth]{Solution/AMC_12_12B_2008_25_Solution_Figure_2.JPG}
                    \end{figure}
                    

 Let $AP$ and $BQ$ meet $CD$ at $X$ and $Y$, respectively.

 Since $\angle APD=90^{\circ}$, $\angle ADP=\angle XDP$, and they share $DP$, triangles $APD$ and $XPD$ are congruent.

 By the same reasoning, we also have that triangles $BQC$ and $YQC$ are congruent.

 Hence, we have $[ABQCDP]=[ABYX]+\frac{[ABCD]-[ABYX]}{2}=\frac{[ABCD]+[ABYX]}{2}$.

 If we let the height of the trapezoid be $x$, we have $[ABQCDP]=\frac{\frac{11+19}{2}\cdot x+\frac{11+7}{2}\cdot x}{2}=12x$.

 Thusly, if we find the height of the trapezoid and multiply it by 12, we will be done.

 Let the projections of $A$ and $B$ to $CD$ be $A'$ and $B'$, respectively.

 We have $DA'+CB'=19-11=8$, $DA'=\sqrt{DA^2-AA'^2}=\sqrt{49-x^2}$, and $CB'=\sqrt{CB^2-BB'^2}=\sqrt{25-x^2}$.

 Therefore, $\sqrt{49-x^2}+\sqrt{25-x^2}=8$. Solving this, we easily get that $x=\frac{5\sqrt{3}}{2}$.

 Multiplying this by 12, we find that the area of hexagon $ABQCDP$ is $30\sqrt{3}$, which corresponds to answer choice $\boxed{B}$.

 



\textbf{Solution 3}


                    \begin{figure}[H]
                        \centering
                        \includegraphics[width=0.5\textwidth]{Solution/AMC_12_12B_2008_25_Solution_Figure_3.png}
                    \end{figure}
                    
$P$ is the intersection of the angle bisectors of $\angle A$ and $\angle D$. By definition, angle bisectors are always equidistant from the sides of the angle, so $P$ is equidistant from $\overline{AB}$, $\overline{AD}$, and $\overline{CD}$.  Likewise, point $Q$ is equidistant from $\overline{AB}$, $\overline{BC}$, and $\overline{CD}$.  Because both points $P$ and $Q$ are equidistant from $\overline{AB}$ and $\overline{CD}$ and the distance between $\overline{AB}$ and $\overline{CD}$ is constant, the common distances from each of the points to the mentioned segments is equal for $P$ and $Q$.  Call this distance $x$.

 The distance between a point and a line is the length of the segment perpendicular to the line with one endpoint on the line and the other on the point.  This means the altitude from $P$ to $\overline{AD}$ is $x$, so the area of $\triangle ADP$ is equal to $\frac12\cdot AD\cdot x=\frac72x$.  Similarly, the area of $\triangle BCQ$ is $\frac12\cdot BC\cdot x=\frac52x$.  The altitude of the trapezoid is $2x$, because it is the sum of the distances from either $P$ or $Q$ to $\overline{AB}$ and $\overline{CD}$.  This means the area of trapezoid $ABCD$ is $\frac12(AB+CD)\cdot2x=\frac12(11+19)\cdot2x=30x$.  Now, the area of hexagon $ABQCDP$ is the area of trapezoid $ABCD$, minus the areas of triangles $ADP$ and $BCQ$.  This is $30x-\frac72x-\frac52x=24x$.  Now it remains to find $x$.

 
                    \begin{figure}[H]
                        \centering
                        \includegraphics[width=0.5\textwidth]{Solution/AMC_12_12B_2008_25_Solution_Figure_4.png}
                    \end{figure}
                    

 We let $R$ and $S$ be the feet of the altitudes of $A$ and $B$, respectively, to $\overline{CD}$.  We define $r=RD$ and $s=SC$.  We know that $AB=RS$, so $RS=11$ and $r+s=19-11=8$.  By the Pythagorean Theorem on $\triangle ADR$ and $\triangle BCS$, we get $r^2+(2x)^2=7^2$ and $s^2+(2x)^2=5^2$, respectively.  Subtracting the second equation from the first gives us $r^2-s^2=49-25=24$.  The left hand side of this equation is a difference of squares and factors to $(r+s)(r-s)$.  We know that $r+s=8$, so $r-s=\frac{24}8=3$.  Now we can solve for $r$ by adding the two equations we just got to see that $2r=11$, or $r=\frac{11}2$.

 We now solve for $x$.  We know that $r^2+(2x)^2=49$, so $(2x)^2=49-\left(\frac{11}2\right)^2=\frac{75}4$ and $2x=\frac{5\sqrt3}2$.  We multiply both sides of this equation by $12$ to get $24x=30\sqrt3$.  However, the area of hexagon $ABQCDP$ is $24x$, so the answer is $30\sqrt 3$, or answer choice $\boxed{B}$.

 



\textbf{Solution 4}


                    \begin{figure}[H]
                        \centering
                        \includegraphics[width=0.5\textwidth]{Solution/AMC_12_12B_2008_25_Solution_Figure_5.png}
                    \end{figure}
                    

 Let $W, X, Y, Z = \overline{AP} \cap \overline{CD}, \overline{BQ} \cap \overline{CD}, \overline{CQ} \cap \overline{AB}, \overline{DP} \cap \overline{AB}$ respectively. Since $\angle{DCQ} = \angle{BCQ}, \angle{CBQ} = \angle{ABQ}$ we have $\angle{QCB} + \angle{CBQ} = 90 \iff \overline{BX} \perp \overline{CY};$ similarly we get $\overline{AW} \perp \overline{DZ}.$ Thus, $\overline{BQ}$ is both an angle bisector and altitude of $\triangle{CBY}$ so $BC = BY.$ Using the same logic on $\triangle{BCX}$ gives $BC = BX \iff BYXC$ is a rhombus; similarly, $ADWZ$ is a rhombus. Then, $[ABQCDP] = [ABCD] - \frac{1}{4}\left([BYXC] + [ADWZ]\right) = 15h - \frac{1}{4}(7h + 12h) = 12h$ where $h$ is the height of trapezoid $ABCD.$ Finding $h$ is the same as finding the altitude to the side of length $8$ in a $5-7-8$ triangle, and using Heron's, the area of such a triangle is $\sqrt{10(5)(3)(2)} = 10 \sqrt{3} = 4h \iff h = \frac{5\sqrt{3}}{2}.$ Multiply to get our answer is $\boxed{\mathrm{B}}.$




\textbf{Solution 5}

Like above solutions, find out the height of $ABCD$ is $\frac{5 \sqrt{3}}{2}.$ Let $M$ be the intersection of the line through $Q$ and parallel to $AB,$ and $N$ the intersection of the line through $P$ and parallel to $AB.$ Angle chasing shows that $M$ is the midpoint of $BC$ and $N$ midpoint of $AD.$ Then from midline theorem,  $M, Q, N$ are collinear, and likewise for $N, P, M.$ Thus, the line through $PQ$ is in fact the midline of $ABCD.$

 
Let $BQ \cap CD = X, AP \cap CD = Y.$ Then, angle chasing shows that $CQ$ not only bisects $BX,$ but is also perpendicular to it. This makes it a perpendicular bisector. The same is true for $DP$ and $AY.$ Thus, $CX = CB = 5,$ and $DY = DA = 7.$ This means $XY = 19 - 5 - 7 = 7.$ We can now find $PQ$ as the midline of $ABXY.$ Thus, $PQ = \frac{1}{2} (11+7) = 9.$

 Now, the answer is simply finding the area of $ABQP$ plus area of $CDPQ.$ This is $\frac{1}{2} \cdot \frac{5 \sqrt{3}}{2} \cdot \frac{11+9}{2} + \frac{1}{2} \cdot \frac{5 \sqrt{3}}{2} \cdot \frac{19+9}{2}  = \boxed{30 \sqrt{3}}.$




\textbf{Solution 6}

Observe that if we reduce the lengths of the parallel sides by the length $PQ$, $P$ and $Q$ will coincide because $PQ \parallel AB.$ When they coincide, they happen to be the incenter of $ABCD$ because all four angle bisectors intersect there. 

 Let the shortened trapezoid be $AB'C'D.$

 To determine the length of $PQ,$ we can use the two tangent theorem to get $AB' + C'D = AD + B'C',$ or $7+5=19-PQ+11-PQ,$ so $PQ = 9.$

 The distance from $P$ to both $AB$ and $CD$ is the same because it is the incenter. To find the height $r$ of both trapezoids, translate $BC$ to $B''C''$ such that $C''$ maps to $D.$ By using Heron's formula and the known base of $19-11=8,$ we get that the altitude of $ABCD$ is $\frac{5\sqrt{3}}{2}.$ Therefore each of the altitudes of $ABQP$ and $CDOQ$ are $\frac{5\sqrt{3}}{4}.$

 Therefore, the area of the hexagon is the sum of the area of the two trapezoids, which is just 
 \[\frac{5\sqrt{3}}{4} \left( \dfrac{19+9}{2}\right) + \frac{5\sqrt{3}}{4} \left( \dfrac{11+9}{2}\right) = \boxed{\text{(B) }30\sqrt{3}}.\]

 - \href{https://artofproblemsolving.com/wiki/index.php/User:Spectraldragon8}{spectraldragon8}

 


\end{enumerate}